% ------------------------------------------------------------------------------------
%                       PROBLEMI DI MODELLAZIONE - ESERCIZI
% ------------------------------------------------------------------------------------
\chapter{Modellazione - Esercizi}

\section{CUP}

Immaginiamo di avere un \textbf{insieme di quartieri} $l\in L = \{1,\ldots , n\}$ ed 
una \textbf{matrice} $L\times L$ che raccoglie i \textbf{tempi di percorrenza} 
(\textit{previsti}) necessari per \textbf{raggiungere un CUP} 
posto in un altro quartiere, come mostrato sotto:
\[
    \begin{matrix}
        & | & 1 & 2 & 3 & 4 & 5 & 6 \\ 
      - & | & - & - & - & - & - & - \\ 
      1 & | & 0 & 10 & 20 & 30 & 30 & 20 \\ 
      2 & | & 10 & 0 & 25 & 33 & 20 & 10 \\ 
      3 & | & \vdots & \vdots & \vdots & \vdots & \vdots & \vdots \\ 
      4 & | & \vdots & \vdots & \vdots & \vdots & \vdots & \vdots \\ 
      5 & | & \vdots & \vdots & \vdots & \vdots & \vdots & \vdots \\ 
      6 & | & \vdots & \vdots & \vdots & \vdots & \vdots & \vdots \\ 
    \end{matrix}
\]

\begin{remark}
    La matrice sarà \textbf{simmetrica}: andare \textit{da} $i$ a $j$ è
    \textit{equivalente} ad andare da $j$ ad $i$.
\end{remark}

E chiamiamo $d_{ij}$ le \textit{entrate della matrice}, ovvero il \textbf{tempo di percorrenza}
per \textbf{raggiunere un CUP} posto in $i$ \textbf{dal quartiere} $j$.
L'obiettivo del problem è:

\begin{exa}[PROBLEMA]
    Vogliamo \textbf{posizionare i CUP nei quartieri} in modo che
    il \textbf{tempo di percorrenza massimo} $T$ da \textit{qualsiasi quartiere}
    sia di \textbf{15 minuti}.
\end{exa}

Per come è descritto il problema:

\begin{remark}
    Notiamo che il problema è essenzialmente un problema
    di \textbf{set covering}: si tratta infatti di \textbf{soddisfare} 
    dei \textbf{vincoli di copertura}.
\end{remark}

Per modellare il problema, cominciamo definendo le \textbf{variabili binarie}:
\[
    x_i = \begin{cases}
        1 & \text{se posizioniamo un CUP in posizione } $i$ \\ 
        0 & \text{altrimenti}
    \end{cases}
\]

Abbiamo quindi:
\[
    \begin{cases}
        \displaystyle\text{min } \sum_{i\in I} x_i \\ 
        \displaystyle\sum_{i\in L : d_{ij} \le T} x_i \ge 1 & \forall j \in L \\
        x_i \in \{0,1\} & \forall i \in L
    \end{cases}
\]

Come nel caso di \textbf{set covering}, vogliamo usare il \textbf{minor numero di CUP}
possibile (equivalente ad usare il minor numero di sottoinsiemi).
Detto ciò, diamo una spiegazione del vincolo imposto:

\begin{remark}
    Per \textbf{ogni quartiere} $j\in L$ deve \textbf{esistere almeno un CUP}
    posto in \textit{un qualche} $i\in L$ tale che il \textbf{tempo di percorrenza}
    sia \textbf{inferiore alla soglia} richiesta.
\end{remark}
